% Auto-generated from 42-Institutional-Replacement.md
% POV: Aisen | Landscape: Dungeon | Location: The Hidden Passage

With Aisen now in civilian status, he could operate with greater freedom from military regulations. But civilian operations in Covenant-controlled territories required navigating different kinds of restrictions: administrative oversight, civilian security monitoring, institutional authority that came through educational and economic structures rather than military command.

The first institutional replacement operations in Southeastern Territory focused on replacing provincial administrators with Covenant-approved personnel. This involved installing Covenant-trained governors, replacing local judiciary with Covenant-certified judges, restructuring education system to align with Covenant institutional priorities.

Provincial resistance networks adapted quickly. They couldn't prevent institutional replacement—they lacked the political authority to do that. But they could position civil society organizations to eventually challenge Covenant institutional control. They could establish alternative educational structures. They could develop parallel administrative structures that would become visible after Covenant institutional control became sufficiently normalized that security monitoring reduced.

Aisen coordinated this work for Kael's resistance network while maintaining civilian cover as what appeared to be an educational coordinator for provincial schools. His position gave him access to educational administration and allowed him to quietly position resistance network members in institutional positions where they could eventually influence how education functioned.

"We're not trying to overthrow Covenant institutional authority," Aisen said to new resistance network members learning institutional subversion approaches. "We're learning how institutions actually function and positioning ourselves to influence that function. When energy becomes available—when security attention focuses elsewhere, when institutional contradictions become visible, when population support for alternatives becomes apparent—then we have infrastructure positioned to facilitate institutional change."

The approach was long-term and patient in ways that military resistance operations had never been. It was also, Aisen realized, probably more strategically effective than armed resistance would have been. Institutions maintained themselves through thousands of small daily decisions made by ordinary people. If resistance networks could influence those decisions, institutional alternatives could gradually become possible.
